\section*{Preface}
\phantomsection
\addcontentsline{toc}{section}{Preface}
\label{sec:preface}

These notes ask how electrons confined in a solid can store, manipulate,
and reveal quantum information. We begin with the physical system that
holds the electrons, determine which of their states can be used, and
then develop the interactions and measurements needed to operate on
those states. The main text gives the results and conditions needed for
this explanation; the appendices develop their derivations and the
more specialized calculations.

\begingroup
\renewcommand{\theequation}{Pref.\arabic{equation}}

\paragraph{Spin and Pauli operators.}
Spin states are labelled by the total-spin quantum number $S$ and
the magnetic quantum number $m$. For a physical electron spin, we
use the ordered basis $(\ket{\uparrow},\ket{\downarrow})$ and
Pauli operators $\sigma_x,\sigma_y,\sigma_z$, with
$\sigma_z\ket{\uparrow}=\ket{\uparrow}$ and
$\sigma_z\ket{\downarrow}=-\ket{\downarrow}$.
The dimensionless spin operator is
$\mathbf s=\boldsymbol\sigma/2$; its angular momentum is
$\hbar\mathbf s$. Spin and Pauli operators are written without hats.

For a logical qubit, the Pauli operators $X,Y,Z$ refer to the ordered
basis $(\ket0,\ket1)$:
\begin{equation}
 X=\begin{pmatrix}0&1\\1&0\end{pmatrix},\qquad
 Y=\begin{pmatrix}0&-i\\i&0\end{pmatrix},\qquad
 Z=\begin{pmatrix}1&0\\0&-1\end{pmatrix}.
 \label{eq:elementary-logical-gate-matrices}
\end{equation}
The physical-spin operators $\sigma_x,\sigma_y,\sigma_z$ have these
same matrices in the physical-spin basis above. Their relation to
$X,Y,Z$ depends on the encoding, whose logical states are defined
in the main text. Subscripts identify the spin or logical qubit
on which an operator acts; identity operators on other factors
are implicit. Two-qubit matrices use the basis order
$(\ket{00},\ket{01},\ket{10},\ket{11})$ unless stated otherwise.

\paragraph{Logical rotations.}
We use $R_k(\vartheta)=\exp(-i\vartheta K/2)$, with
$(k,K)=(x,X),(y,Y),(z,Z)$:
\begin{equation}
 \begin{aligned}
 R_x(\vartheta)&=
 \begin{pmatrix}
 \cos(\vartheta/2)&-i\sin(\vartheta/2)\\
 -i\sin(\vartheta/2)&\cos(\vartheta/2)
 \end{pmatrix},\\[3pt]
 R_y(\vartheta)&=
 \begin{pmatrix}
 \cos(\vartheta/2)&-\sin(\vartheta/2)\\
 \sin(\vartheta/2)&\cos(\vartheta/2)
 \end{pmatrix},\\[3pt]
 R_z(\vartheta)&=
 \begin{pmatrix}e^{-i\vartheta/2}&0\\0&e^{i\vartheta/2}\end{pmatrix}.
 \end{aligned}
 \label{eq:logical-rotation-matrices}
\end{equation}
Encoding subscripts identify the logical states on which these
operators act. Products act from right to left.

\paragraph{Evolution and phases.}
For a Hamiltonian $H(t)$ applied over an interval $\tau$, the
evolution operator is
\begin{equation}
 U(\tau)=\mathcal T\exp\!\left[-\frac{i}{\hbar}
                    \int_0^\tau H(t)\,dt\right].
 \label{eq:gate-time-evolution}
\end{equation}
The time-ordering operator $\mathcal T$ places later-time operators
to the left.
When specifying gates on an isolated computational space,
an overall phase does not change their physical action. A scalar
term $c(t)I$ in the Hamiltonian contributes only the phase
$e^{-i\int_0^\tau c(t)\,dt/\hbar}$ and may therefore be omitted.
Relative phases between computational basis states are part of
the gate and must be retained.

\paragraph{Pauli and Hadamard gates.}
The operators $X,Y,Z$ also define the Pauli gates. The Hadamard gate is
\begin{equation}
 H_{\mathrm H}=\frac{1}{\sqrt2}
                  \begin{pmatrix}1&1\\1&-1\end{pmatrix}.
 \label{eq:hadamard-gate-matrix}
\end{equation}
Their decompositions into rotations are
\begin{equation}
 X=iR_x(\pi),\qquad Y=iR_y(\pi),\qquad Z=iR_z(\pi),\qquad
 H_{\mathrm H}=iR_y(\pi/2)R_z(\pi).
 \label{eq:elementary-single-qubit-gates}
\end{equation}
More generally, the Euler product
$R_z(\vartheta_1)R_y(\vartheta_2)R_z(\vartheta_3)$ gives an
arbitrary single-qubit unitary up to global phase.

\paragraph{SWAP and \texorpdfstring{$\sqrt{\mathrm{SWAP}}$}{square-root SWAP}.}
SWAP exchanges the two logical states,
$\mathrm{SWAP}\ket{ab}=\ket{ba}$ for $a,b\in\{0,1\}$.
We use the following choice of its square root:
\begin{equation}
 \mathrm{SWAP}=\begin{pmatrix}
  1&0&0&0\\0&0&1&0\\0&1&0&0\\0&0&0&1
 \end{pmatrix},\qquad
 \sqrt{\mathrm{SWAP}}=\frac12\begin{pmatrix}
  2&0&0&0\\0&1+i&1-i&0\\0&1-i&1+i&0\\0&0&0&2
 \end{pmatrix}.
 \label{eq:exchange-gate-matrices}
\end{equation}

\paragraph{CZ and CNOT gates.}
CZ adds a minus sign to $\ket{11}$. CNOT flips the target qubit
when the control is $\ket1$. With the first qubit as control,
\begin{equation}
 \mathrm{CZ}=\begin{pmatrix}
  1&0&0&0\\0&1&0&0\\0&0&1&0\\0&0&0&-1
 \end{pmatrix},\qquad
 \mathrm{CNOT}_{1\to2}=\begin{pmatrix}
  1&0&0&0\\0&1&0&0\\0&0&0&1\\0&0&1&0
 \end{pmatrix}.
 \label{eq:controlled-logical-gate-matrices}
\end{equation}
For logical qubits $A,B$, the same CNOT is
\begin{equation}
 \mathrm{CNOT}_{A\to B}
   =(\ket0\bra0)_A\otimes I_B
    +(\ket1\bra1)_A\otimes X_B.
 \label{eq:exchange-only-cnot-definition}
\end{equation}
It maps $\ket{ab}$ to $\ket{a,b\mathbin\oplus a}$, where
$\oplus$ denotes addition modulo two. Hadamard gates on the target
convert CZ into CNOT:
\begin{equation}
 \mathrm{CNOT}_{A\to B}
   =(I_A\otimes H_{\mathrm H,B})\,\mathrm{CZ}_{AB}\,
     (I_A\otimes H_{\mathrm H,B}).
 \label{eq:exchange-cnot-construction}
\end{equation}

\endgroup
